% ===================== AUDIT POKRYTÍ =====================
\section{Audit pokrytí požadavků}
\label{sec:audit}

\begin{infobox}[title={Závěr auditu}]
Prošel jsem \textbf{všechny} body obou přípravných dokumentů (oficiální požadavky SZZ, verze 28.\,7.\,2025) a porovnal je s obsahem poznámek notes-ipp. \textbf{Každé požadované téma je v poznámkách obsažené.} Jediné formálně chybějící je \emph{Hilbertovský kalkul}, který je ovšem v požadavcích uvedený jen jako \emph{alternativa} k tablu/rezoluci (a ty obě pokryté jsou) - viz \ref{sec:vyjimka}. Šest míst poznámky vědomě nechávají na vlastní procvičení (výpočetní dovednosti) - viz \ref{sec:procvic}.
\end{infobox}

\subsection{Metodika}
Pro každý okruh jsem vzal seznam podtémat z oficiálních požadavků a ověřil, že odpovídající látka v poznámkách existuje (strukturálně i obsahově - členění poznámek kopíruje oficiální osnovu). U níže uvedených seznamů znamená \checkmark{}, že téma je v textu zpracované.

\subsection{Souhrn}
\begin{center}
\begin{tabular}{@{}lcc@{}}
\toprule
\textbf{Okruh} & \textbf{Témat v požadavcích} & \textbf{Pokryto} \\
\midrule
I/1 Diferenciální a integrální počet & 8 & \checkmark{} vše \\
I/2 Algebra a lineární algebra & 9 & \checkmark{} vše \\
I/3 Diskrétní matematika & 7 & \checkmark{} vše \\
I/4 Teorie grafů & 9 & \checkmark{} vše \\
I/5 Pravděpodobnost a statistika & 8 & \checkmark{} vše \\
I/6 Logika & 6 & \checkmark{} vše* \\
II/1 Automaty a jazyky & 4 & \checkmark{} vše \\
II/2 Algoritmy a datové struktury & 6 & \checkmark{} vše \\
II/3 Programovací jazyky & 8 & \checkmark{} vše \\
II/4 Architektura počítačů a OS & 7 & \checkmark{} vše \\
III\phantom{/0} Základy umělé inteligence & 9 & \checkmark{} vše \\
IV\phantom{/0} Strojové učení & 9 & \checkmark{} vše \\
\bottomrule
\end{tabular}
\end{center}
{\small *U logiky chybí jen Hilbertovský kalkul, uvedený v požadavcích jako alternativa (tablo a rezoluce pokryté jsou).}

\subsection{Detailní kontrola po tématech}

\textbf{Část I - Společná matematika}
\begin{multicols}{2}\small
\textit{1. Diferenciální a integrální počet}
\begin{itemize}[label=\checkmark,leftmargin=1.3em,itemsep=0pt]
  \item posloupnosti a limity, aritmetika limit
  \item věta o dvou policajtech
  \item řady; geometrická a harmonická řada
  \item limita funkce, spojitost (mezihodnoty, maximum)
  \item derivace, l'Hospital, průběh funkce
  \item Taylorův polynom (limitní forma)
  \item primitivní funkce, Riemann/Newton
  \item aplikace integrálu (obsah, objem, povrch, délka)
\end{itemize}
\textit{2. Algebra a lineární algebra}
\begin{itemize}[label=\checkmark,leftmargin=1.3em,itemsep=0pt]
  \item grupy, permutace, tělesa (konečná)
  \item soustavy rovnic, Gauss(-Jordan)
  \item matice, hodnost, inverze
  \item vektorové prostory, Steinitz, báze, dimenze
  \item lineární zobrazení, obraz/jádro, isomorfismus
  \item skalární součin, Gram-Schmidt, projekce
  \item determinanty (Laplace, geometrie)
  \item vlastní čísla, diagonalizace, spektrální rozklad
  \item pozitivně (semi)definitní, Cholesky
\end{itemize}
\textit{3. Diskrétní matematika}
\begin{itemize}[label=\checkmark,leftmargin=1.3em,itemsep=0pt]
  \item vlastnosti relací
  \item ekvivalence a rozklady
  \item částečná uspořádání, výška/šířka
  \item funkce a jejich počty
  \item kombinační čísla, binomická věta
  \item inkluze a exkluze (i aplikace)
  \item Hallova věta a SRR
\end{itemize}
\textit{4. Teorie grafů}
\begin{itemize}[label=\checkmark,leftmargin=1.3em,itemsep=0pt]
  \item základní pojmy a příklady grafů
  \item souvislost, komponenty, vzdálenost
  \item stromy a charakteristiky
  \item rovinné grafy, Eulerova formule
  \item barevnost a klikovost
  \item Mengerova věta (obě verze)
  \item orientované grafy, silná/slabá souvislost
  \item toky, Ford-Fulkerson
\end{itemize}
\textit{5. Pravděpodobnost a statistika}
\begin{itemize}[label=\checkmark,leftmargin=1.3em,itemsep=0pt]
  \item prostor, podmíněná pst, Bayes
  \item náhodné veličiny, distribuce, hustota
  \item střední hodnota, Markovova nerovnost
  \item rozptyl (i pro součet)
  \item konkrétní rozdělení
  \item LLN a centrální limitní věta
  \item bodové a intervalové odhady
  \item testování hypotéz, chyby I/II, hladina
\end{itemize}
\textit{6. Logika}
\begin{itemize}[label=\checkmark,leftmargin=1.3em,itemsep=0pt]
  \item syntaxe, CNF/DNF/PNF, SAT/rezoluce
  \item sémantika, splnitelnost, důsledek
  \item extenze, konzervativnost, skolemizace
  \item dokazatelnost: tablo, rezoluce
  \item kompaktnost a úplnost
  \item rozhodnutelnost
\end{itemize}
\end{multicols}

\textbf{Část II - Společná informatika}
\begin{multicols}{2}\small
\textit{1. Automaty a jazyky}
\begin{itemize}[label=\checkmark,leftmargin=1.3em,itemsep=0pt]
  \item regulární jazyky (gramatiky, DFA/NFA, regexy)
  \item bezkontextové (gramatiky, zásobníkový automat)
  \item rekurzivně spočetné (typ 0, Turing, nerozhodnutelnost)
  \item Chomského hierarchie
\end{itemize}
\textit{2. Algoritmy a datové struktury}
\begin{itemize}[label=\checkmark,leftmargin=1.3em,itemsep=0pt]
  \item složitost, asymptotická notace
  \item P, NP, převoditelnost, NP-úplnost
  \item rozděl a panuj, Master theorem
  \item vyhledávací stromy, AVL
  \item třídění a dolní odhad
  \item grafové algoritmy (BFS/DFS, Dijkstra, MST, toky)
\end{itemize}
\textit{3. Programovací jazyky}
\begin{itemize}[label=\checkmark,leftmargin=1.3em,itemsep=0pt]
  \item abstrakce, dědičnost, polymorfismus
  \item typy, reference, boxing
  \item generika a funkcionální prvky
  \item zdroje a výjimky
  \item životní cyklus, GC, reference counting
  \item vlákna a synchronizace
  \item vtable a implementace OO
  \item bytecode, JIT/AOT, linkování, knihovny
\end{itemize}
\textit{4. Architektura počítačů a OS}
\begin{itemize}[label=\checkmark,leftmargin=1.3em,itemsep=0pt]
  \item architektura, reprezentace čísel, adresování
  \item instrukční sada a vazba na vyšší jazyk
  \item privilegovaný režim, jádro
  \item periferie: PIO, přerušení
  \item procesy/vlákna, kontext, plánování
  \item virtuální paměť, stránkování, ochrana
  \item soubory, race conditions, zámky
\end{itemize}
\end{multicols}

\textbf{Část III - Základy UI \quad / \quad Část IV - Strojové učení}
\begin{multicols}{2}\small
\textit{III. Základy umělé inteligence}
\begin{itemize}[label=\checkmark,leftmargin=1.3em,itemsep=0pt]
  \item prohledávání (DFS/BFS/UCS, A*, heuristiky)
  \item CSP, AC-3, MAC
  \item logické uvažování (DPLL, řetězení, rezoluce)
  \item pravděpodobnostní uvažování, Bayesovy sítě
  \item reprezentace znalostí, HMM vs.\ DBN
  \item automatické plánování
  \item MDP, Bellman, value/policy iteration, POMDP
  \item hry (minimax, alfa-beta), Nash, mechanism design
  \item ML přehledově (učení, stromy, SVM, EM, RL)
\end{itemize}
\textit{IV. Strojové učení}
\begin{itemize}[label=\checkmark,leftmargin=1.3em,itemsep=0pt]
  \item učení s učitelem, metriky, regularizace
  \item k-nejbližších sousedů
  \item lineární regrese (LSQ, SGD)
  \item logistická regrese (sigmoid, softmax)
  \item rozhodovací stromy, prořezávání
  \item SVM (separabilní/neseparabilní, jádra, multiclass)
  \item kombinace modelů (bagging, boosting, random forest)
  \item statistické testy (t-test, chí-kvadrát)
  \item učení bez učitele (k-means, hierarchické, PCA)
\end{itemize}
\end{multicols}

\subsection{Místa, která poznámky nechávají na vlastní procvičení}
\label{sec:procvic}
Nejde o chybějící definice, ale o \textbf{výpočetní dovednosti}, které text vědomě nerozepisuje. V příslušných okruzích jsou navíc označené žlutým rámečkem přímo u tématu.
\begin{enumerate}[leftmargin=1.5em]
  \item Derivační pravidla (součin, podíl, složená funkce) - I/1.
  \item Integrační techniky: substituce a per partes - I/1.
  \item Odhad součtu řady přes integrál (konkrétní postup) - I/1.
  \item Vlastní výpočet determinantu - I/2.
  \item Překlad konstrukcí vyššího jazyka do instrukcí procesoru - II/4.
  \item Opačný směr: poznat konstrukci z dané sekvence instrukcí - II/4.
\end{enumerate}

\subsection{Jediná formální výjimka: Hilbertovský kalkul}
\label{sec:vyjimka}
\begin{gapbox}
Požadavky u okruhu \emph{Logika} zmiňují dokazovací systémy „tablo, rezoluce \textbf{nebo} Hilbertovský kalkul“. Poznámky podrobně rozebírají \textbf{tablo metodu i rezoluci}, čímž je požadavek splněn. \textbf{Hilbertovský (axiomatický) kalkul} v poznámkách \textbf{není}. Protože je uvedený jen jako alternativa, na zkoušku to nevadí; pokud bys ho přesto chtěl umět, doplň si: jazyk s několika \emph{axiomovými schématy} (implikační fragment) a jediným odvozovacím pravidlem \emph{modus ponens}; důkaz je posloupnost formulí, kde každá je buď axiom, předpoklad, nebo plyne z předchozích pravidlem; věta o dedukci. Toto si \textbf{ověř} v oficiálních materiálech NAIL062, pokud bys na něj chtěl spoléhat.
\end{gapbox}
